\section{Descripción del Problema de Investigación}
\noindent Poder estudiar las inestabilidades en plasmas astrofísicos mediante simulaciones numéricas se ha convertido en el método preferido en los últimos años. Esto facilita la comprensión de diferentes fenómenos que ocurren en objetos tan densos como las estrellas de neutrones y los discos de acreción, y ayuda a explicar los chorros relativistas.

Un ejemplo bien conocido es la inestabilidad de Kelvin-Helmholtz (KHI). Esta surge cuando hay grandes diferencias o saltos de velocidad dentro de un fluido o plasma magnetizado. Tal inestabilidad llevará a la mezcla turbulenta e incluso incrementará la intensidad de los campos magnéticos. Entonces, no solo puede conducir a aumentos mayores en las intensidades del campo magnético, sino que también le otorga una naturaleza muy caótica, o de “viento” turbulento, a tubos de flujo originalmente bien comportados.

En el trabajo de \citep{osmanov-2008} se menciona que cuando los plasmas son tanto magnetizados como relativistas, como aquellos que se encuentran cerca de púlsares o en eyecciones de masa a velocidades muy altas, se necesitan modelos más complejos para describir con precisión su comportamiento. Uno de esos modelos es la Magnetohidrodinámica Resistiva Relativista (RRMHD), que toma en consideración el efecto de la resistividad finita en fenómenos no ideales como la reconexión magnética.

Investigaciones como la de \citep{pimentel-2019} han mostrado que la resistividad juega un papel importante en la evolución de la inestabilidad KH en estos plasmas relativistas. La resistividad facilita la reconexión entre áreas de esfuerzo de corte de velocidad paralelas a los campos magnéticos e impacta directamente la tasa a la cual las perturbaciones crecen. Los autores observaron, gracias a simulaciones bidimensionales de alta resolución, que la incorporación de resistividad explícita altera fundamentalmente tanto la forma como la evolución no lineal de la inestabilidad. Tales hallazgos son muy importantes porque podrían impactar cómo los modelos que explican el transporte de energía y momento funcionan para chorros relativistas y ciertos tipos de estrellas con magnetósferas.

Además, otros trabajos como el de \citep{mizuno-2013} han avanzado sustancialmente en el estudio de la inestabilidad KH en flujos relativistas que también están magnetizados. Su conclusión es que la dirección del campo magnético y la diferencia de densidad entre las capas de corte tienen un efecto muy importante en la estabilidad del sistema. También afectan grandemente cómo el sistema transita a un estado turbulento. Estas simulaciones han proporcionado la base para muchos de los modelos de astrofísica computacional actuales.

En el trabajo de \citep{bucciantini-2006} realizaron un estudio numérico detallado de la inestabilidad KH en chorros magnetizados, considerando diferentes configuraciones del campo magnético. Descubrieron que campos magnéticos orientados paralelamente al flujo pueden estabilizar algunos modos, mientras que componentes perpendiculares permiten el desarrollo más libre de inestabilidades. Su trabajo destaca cómo la geometría del campo magnético es determinante para la evolución del sistema y refuerza la necesidad de incluir efectos magnetohidrodinámicos completos en las simulaciones .

Por otro lado, el trabajo de \citep{malagoli-1996} quienes realizaron un estudio lineal detallado de la estabilidad de chorros relativistas. Trataron varios perfiles de velocidad, densidad y presión. Su investigación reveló que la dependencia de la inestabilidad en la relación entre el chorro y el medio externo es fuerte. También identificaron escalas críticas para la transición a la región no lineal y mostraron cómo configuraciones particulares pueden incluso amplificar modos helicoidales o transversales. Esto ha sido crucial para la interpretación de observaciones radioastronómicas de características en chorros y para que cálculos tridimensionales posteriores sean válidos.

Para modelar de forma más realista los plasmas astrofísicos, es necesario ir más allá de la Magnetohidrodinámica (MHD) ideal, que asume que el fluido es un conductor perfecto. Sin embargo, los modelos de MHD resistiva presentan un gran desafío numérico: las ecuaciones adquieren términos rígidos (stiff) que describen efectos difusivos en escalas de tiempo muy cortas, lo que complica enormemente la evolución temporal en las simulaciones. Estos términos pueden dominar sobre la parte puramente hiperbólica y exigen pasos de tiempo extremadamente pequeños para mantener la estabilidad, haciendo que los métodos explícitos sean impracticables.

El trabajo de \citep{palenzuela-2009} propone una solución a este problema. Los autores introducen un novedoso enfoque para la solución numérica de las ecuaciones de la MHD resistiva relativista basado en los métodos Runge-Kutta implícitos-explícitos (IMEX). Este esquema trata la parte hidrodinámica de las ecuaciones con un método explícito y los términos de relajación rígidos con uno implícito , evitando las severas restricciones en el paso de tiempo. Una de sus contribuciones más importantes es que este enfoque permite tratar, dentro de un marco unificado, regiones dominadas tanto por la presión del fluido (como el interior de objetos compactos) como por la presión magnética (como sus magnetosferas). Mediante una serie de pruebas numéricas, demostraron la robustez de su implementación y destacaron su eficiencia frente a técnicas alternativas como la descomposición de Strang, la cual se mostró inestable para flujos discontinuos con altas conductividades.

